\par\noindent\textbf{(1)} 定义条件平移 \(S|x,\uparrow\rangle=|x+1,\uparrow\rangle\)、\(S|x,\downarrow\rangle=|x-1,\downarrow\rangle\)。在Bloch波数 \(k\) 上，\(S(k)=\operatorname{diag}(e^{ik},e^{-ik})\)，一步算符 \(U(k)=S(k)H\)；因为两因子均酉，\(U(k)^\dagger U(k)=I\)。
\par\noindent\textbf{(2)} 矩阵行列式 \(\det U(k)=-1\)，迹 \(\operatorname{tr}U(k)=i\sqrt2\sin k\)。将本征值写 \(e^{i\omega(k)}\) 与 \(-e^{-i\omega(k)}\)，由迹得 \(2i\sin\omega=i\sqrt2\sin k\)，故 \(\sin\omega(k)=\sin k/\sqrt2\)，群速 \(d\omega/dk=\cos k/\sqrt{1+\cos^2 k}\)。
\par\noindent\textbf{(3)} 由 \(|v_g(k)|=|\cos k|/\sqrt{1+\cos^2k}\)，最大值 \(v_{\max}=1/\sqrt2\)，在 \(k=0,\pi\) 等位置达到。经过 \(t\) 步，具有宽频Bloch成分的波包传播支持主要集中在 \(|x|\lesssim t/\sqrt2\) 内，宽度随时间线性增长而非经典无偏随机行走的 \(\sqrt t\)。
\par\noindent\textbf{(4)} 经典每步独立向左或右移动的位移方差为 \(\operatorname{Var}(x)=t\)，典型宽度 \(\sqrt t\)。量子行走因不同路径振幅相干叠加而呈 \(O(t)\) 弹道展宽；若每步硬币完全退相干，干涉项消失并逐渐恢复扩散标度，不能把酉行走的量子优势自动推广到含强噪声的实验。
